\UnnumberedChapter{LỜI MỞ ĐẦU}

Trong bối cảnh ngành Xây dựng Việt Nam đang đẩy mạnh quá trình chuyển đổi số, việc ứng dụng công nghệ thông tin vào các nghiệp vụ truyền thống là một xu thế tất yếu. Công tác lập dự toán và quản lý chi phí vật tư cơ điện (MEP) hiện nay vẫn phụ thuộc chủ yếu vào thao tác thủ công của kỹ sư: tra cứu hàng nghìn file báo giá Excel, đối chiếu chéo các thông số kỹ thuật và ước lượng giá dựa trên kinh nghiệm cá nhân. Quy trình này không chỉ tốn nhiều thời gian mà còn tiềm ẩn rủi ro sai sót, ảnh hưởng trực tiếp đến năng lực cạnh tranh của doanh nghiệp khi tham gia đấu thầu.

Sự phát triển vượt bậc của công nghệ Trí tuệ Nhân tạo (AI), đặc biệt là các Mô hình Ngôn ngữ Lớn (LLM) như Google Gemini và Google Gemma, đã mở ra cơ hội áp dụng phương pháp Retrieval-Augmented Generation (RAG) để giải quyết bài toán tra cứu và tư vấn giá vật tư một cách tự động, nhanh chóng và có cơ sở dữ liệu minh bạch.

Được sự phân công của khoa Công nghệ Thông tin và sự hướng dẫn trực tiếp của ThS. Lê Thị Thùy Trang, nhóm sinh viên thực hiện đề tài: \textit{``Ứng dụng RAG và LLM trong tự động hóa tư vấn dự toán vật tư cơ điện''}. Báo cáo trình bày quá trình nghiên cứu lý thuyết, phân tích thiết kế, cài đặt hệ thống và đánh giá hiệu năng thực tế của sản phẩm phần mềm Price Advisor AI trên bộ dữ liệu 500 mẫu vật tư MEP thực tế.

Nội dung báo cáo được tổ chức thành 4 chương chính:
\begin{itemize}
    \item \textbf{Chương 1:} Tổng quan bài toán -- trình bày bối cảnh, mục tiêu và phạm vi nghiên cứu.
    \item \textbf{Chương 2:} Cơ sở lý thuyết và công nghệ -- phân tích các nền tảng lý thuyết về LLM, RAG, Dense Retrieval và MLOps.
    \item \textbf{Chương 3:} Phân tích, thiết kế và cài đặt hệ thống -- mô tả kiến trúc tổng thể, các thuật toán cốt lõi và giao diện người dùng.
    \item \textbf{Chương 4:} Thử nghiệm và đánh giá -- trình bày kịch bản Benchmark, kết quả định lượng và phân tích lỗi chuyên sâu.
\end{itemize}
